\documentclass[11pt]{article}
\usepackage[margin=1in]{geometry}
\usepackage{amsmath,amssymb}
\usepackage{enumitem}
\usepackage{hyperref}

\title{Vision-Guided Turret IK Plan}
\author{}
\date{}

\begin{document}
\maketitle

\section*{Goal}
Add a vision-guided control mode in which:
\begin{itemize}[leftmargin=1.5em]
  \item Raspberry Pi 2 runs a camera and AprilTag perception stack.
  \item The camera tracks an AprilTag rigidly attached to the spiral zipper end effector.
  \item Pi 2 publishes the measured 3D tag pose.
  \item Raspberry Pi 1 runs closed-loop inverse kinematics (IK) to move the turret so that the end effector reaches a desired goal position near the camera.
\end{itemize}

\section*{High-Level Architecture}
The system should be split into perception and control:
\begin{itemize}[leftmargin=1.5em]
  \item \textbf{Pi 2 (Perception):} camera driver, camera calibration, AprilTag detection, 3D pose estimation, TF publication.
  \item \textbf{Pi 1 (Control):} receives the current end-effector pose, computes IK, and drives yaw, pitch, and extension through the existing turret controller.
\end{itemize}

This should be implemented as a distributed ROS 2 system. Both Pis share the same ROS domain and exchange messages and TF frames over the network.

\section*{Frame Definitions}
Use TF consistently. Recommended frames:
\begin{itemize}[leftmargin=1.5em]
  \item \texttt{base\_link}: turret base frame.
  \item \texttt{yaw\_link}: frame after yaw rotation.
  \item \texttt{pitch\_link}: frame at turret pitch pivot.
  \item \texttt{end\_effector}: actual spiral zipper tip / tool point.
  \item \texttt{tag\_link}: AprilTag frame rigidly attached to the end effector.
  \item \texttt{camera\_link}: physical camera frame.
  \item \texttt{camera\_optical\_frame}: ROS optical camera frame.
\end{itemize}

Known or measured transforms:
\begin{itemize}[leftmargin=1.5em]
  \item \(T_{\text{base}\_\text{camera}}\): static transform from turret base to camera.
  \item \(T_{\text{tag}\_\text{ee}}\): static transform from the tag center to the actual end-effector tool point.
  \item \(T_{\text{camera}\_\text{tag}}\): measured dynamically from AprilTag detection.
\end{itemize}

Then:
\[
T_{\text{base}\_\text{ee}} =
T_{\text{base}\_\text{camera}}
\; T_{\text{camera}\_\text{tag}}
\; T_{\text{tag}\_\text{ee}}
\]

This gives the measured end-effector pose in the turret base frame.

\section*{Perception Pipeline on Pi 2}
\subsection*{Required outputs}
Pi 2 must publish:
\begin{itemize}[leftmargin=1.5em]
  \item rectified camera image,
  \item camera intrinsics on \texttt{/camera/camera\_info},
  \item tag pose in camera frame,
  \item optionally end-effector pose in base frame if Pi 2 also knows the camera extrinsic.
\end{itemize}

\subsection*{Recommended ROS 2 stack}
\begin{itemize}[leftmargin=1.5em]
  \item \textbf{Camera driver:} \texttt{camera\_ros} with \texttt{libcamera} for Raspberry Pi CSI cameras.
  \item \textbf{Calibration:} \texttt{camera\_calibration}.
  \item \textbf{AprilTag detection:} a ROS 2 AprilTag node that consumes rectified images plus \texttt{CameraInfo} and publishes 3D pose.
  \item \textbf{TF:} \texttt{tf2\_ros} for static and dynamic transforms.
\end{itemize}

\subsection*{Calibration requirements}
Accurate tag pose estimation requires:
\begin{itemize}[leftmargin=1.5em]
  \item camera intrinsics (focal lengths, principal point),
  \item distortion coefficients,
  \item correct tag size in meters,
  \item the static camera extrinsic relative to the turret base or world.
\end{itemize}

\section*{Control Pipeline on Pi 1}
\subsection*{Inputs}
Pi 1 should consume:
\begin{itemize}[leftmargin=1.5em]
  \item measured current end-effector pose \(T_{\text{base}\_\text{ee}}\),
  \item desired end-effector goal pose \(T_{\text{base}\_\text{goal}}\) or desired goal position \(p_{\text{goal}} = [x, y, z]^T\).
\end{itemize}

\subsection*{Outputs}
Pi 1 computes and commands:
\begin{itemize}[leftmargin=1.5em]
  \item yaw angle,
  \item pitch angle,
  \item spiral zipper extension.
\end{itemize}

\subsection*{Closed-loop logic}
Do not use one-shot IK only. Use closed-loop updates:
\begin{enumerate}[leftmargin=1.7em]
  \item Measure current end-effector pose from vision.
  \item Compute goal error in \texttt{base\_link}.
  \item Solve IK to get desired turret joint values.
  \item Command the turret toward those values.
  \item Repeat at \(10\) to \(30\) Hz.
\end{enumerate}

This compensates for model mismatch, cable compliance, gravity, and unmodeled mechanical offsets.

\section*{Forward Kinematics}
The turret has three commanded degrees of freedom:
\begin{itemize}[leftmargin=1.5em]
  \item yaw \(\psi\),
  \item pitch \(\theta\),
  \item extension \(a\).
\end{itemize}

Let the pitch-plane end-effector position relative to the pitch pivot be
\[
p_{\text{plane}} =
\begin{bmatrix}
a \cos \theta \\
a \sin \theta
\end{bmatrix}
\]
with the turret pivot offset included from the base frame.

Then yaw rotates this planar position around the base vertical axis:
\[
p_{\text{ee}} = R_z(\psi)\, p_{\text{plane}} + p_{\text{pivot}}
\]
where \(p_{\text{pivot}}\) is the fixed base-to-pivot offset in 3D.

The first implementation should focus on \textbf{position-only IK}. End-effector orientation can be added later if needed.

\section*{Inverse Kinematics Strategy}
\subsection*{Stage 1: solve yaw analytically}
Given the desired point \(p_{\text{goal}} = [x, y, z]^T\),
\[
\psi = \operatorname{atan2}(y, x)
\]

\subsection*{Stage 2: reduce to the pitch-extension plane}
Project the target into the yaw-aligned 2D plane:
\[
r = \sqrt{x^2 + y^2}
\]
Then solve for \(a\) and \(\theta\) using the turret planar geometry.

\subsection*{Stage 3: planar IK}
Two practical options:
\begin{itemize}[leftmargin=1.5em]
  \item \textbf{Analytic IK:} if the final FK expression is simple and well-conditioned.
  \item \textbf{Numerical IK:} use nonlinear least squares on \((a, \theta)\) with planar FK.
\end{itemize}

For the current system, numerical planar IK is likely the fastest robust path because:
\begin{itemize}[leftmargin=1.5em]
  \item the real mechanism already has offsets and backoff behavior,
  \item measured tag position can be used directly to close the loop,
  \item the numerical method tolerates modest model error better.
\end{itemize}

\subsection*{Recommended first IK objective}
Given measured or desired target position \(p_{\text{goal}}\), solve:
\[
\min_{\psi,\theta,a} \left\| FK(\psi,\theta,a) - p_{\text{goal}} \right\|^2
\]
subject to mechanical limits:
\[
\psi_{\min} \le \psi \le \psi_{\max}, \quad
\theta_{\min} \le \theta \le \theta_{\max}, \quad
a_{\min} \le a \le a_{\max}
\]

\section*{Why Closed-Loop IK Is Preferred}
The real turret is not perfectly modeled:
\begin{itemize}[leftmargin=1.5em]
  \item pitch cable compliance,
  \item friction and stiction,
  \item gravity on the spiral zipper,
  \item backlash and deadband,
  \item camera extrinsic errors,
  \item AprilTag mounting offset.
\end{itemize}

Closed-loop vision feedback turns these into correctable residuals instead of requiring exact open-loop modeling.

\section*{What Must Be Measured}
Before full vision-guided IK will work well, measure:
\begin{itemize}[leftmargin=1.5em]
  \item AprilTag physical edge size in meters,
  \item tag center offset from the true end-effector point,
  \item camera pose relative to turret base,
  \item pitch pivot location relative to \texttt{base\_link},
  \item turret joint limits and safe stand-off to the camera.
\end{itemize}

\section*{Safety and Operating Constraints}
\begin{itemize}[leftmargin=1.5em]
  \item The camera goal position should not be the camera optical center. Use a safe desired stand-off distance.
  \item Enforce minimum and maximum extension, pitch, and yaw.
  \item Reject IK outputs that move through hard limits or singular regions.
  \item Start with low-speed motion and software clamps.
  \item Use filtered vision measurements to avoid jitter from single-frame AprilTag noise.
\end{itemize}

\section*{Implementation Milestones}
\begin{enumerate}[leftmargin=1.7em]
  \item Bring up camera on Pi 2 and verify image stream.
  \item Calibrate camera intrinsics and save a stable calibration file.
  \item Run AprilTag detection and verify \(T_{\text{camera}\_\text{tag}}\).
  \item Define \texttt{tag\_link} to \texttt{end\_effector} offset.
  \item Measure and publish \(T_{\text{base}\_\text{camera}}\).
  \item Publish measured \(T_{\text{base}\_\text{ee}}\).
  \item Implement FK on Pi 1 and compare predicted versus measured end-effector position.
  \item Implement IK for position-only control.
  \item Add slow closed-loop motion to a fixed goal.
  \item Improve filtering, clamping, and convergence behavior.
\end{enumerate}

\section*{Recommended First Working Version}
The first practical version should use:
\begin{itemize}[leftmargin=1.5em]
  \item one fixed camera,
  \item one AprilTag on the end effector,
  \item position-only tracking,
  \item closed-loop IK updates at a modest rate,
  \item existing turret low-level actuation as the inner loop.
\end{itemize}

This gives a direct path to a functional visual servoing system without solving the full orientation-control problem up front.

\end{document}
